% Escala de magnitud aparente — la escala al revés de Hiparco.
%
% Diagrama propio en TikZ (sin licencias de terceros). El truco pedagógico: el
% tamaño del punto DECRECE al avanzar la escala, así se ve de un golpe que
% «más número = menos brillo». Marca además dónde muere el ojo humano en el
% cielo urbano de Cartago frente a una vereda oscura.
%
% Se incrusta vía \guiaDiagrama desde el bloque `recursos:` del YAML.
\begin{tikzpicture}[scale=1,>=stealth,font=\sffamily]

  % x = (magnitud + 2) * 1.15  →  de -2 a 7
  \def\mx#1{(#1+2)*1.15}

  % ── Eje ──────────────────────────────────────────────────────────────
  \draw[thick,->,milcNegro!70] (-0.5,0) -- (11.4,0);
  \node[font=\scriptsize\bfseries,milcNegro!70,right] at (11.4,0) {magnitud};

  \foreach \m in {-2,-1,0,1,2,3,4,5,6,7} {
    \draw[milcNegro!60] ({\mx{\m}},0.13) -- ({\mx{\m}},-0.13);
    \node[font=\scriptsize,milcNegro!70,below] at ({\mx{\m}},-0.15) {$\m$};
  }

  % ── Flechas de sentido ───────────────────────────────────────────────
  \node[font=\scriptsize\bfseries,milcMostaza!70!black,above] at (0.4,0.95)
    {$\leftarrow$ más brillante};
  \node[font=\scriptsize\bfseries,milcNegro!55,above] at (10.2,0.95)
    {más débil $\rightarrow$};

  % ── Astros: el radio cae con la magnitud ─────────────────────────────
  % Sirio (-1,5) · Vega (0) · límite urbano (4) · límite del ojo (6)
  \fill[milcMostaza,draw=milcNegro!50] ({\mx{-1.5}},0.5) circle (5.5pt);
  \node[font=\scriptsize\bfseries,above] at ({\mx{-1.5}},0.72) {Sirio};

  \fill[milcMostaza,draw=milcNegro!50] ({\mx{0}},0.5) circle (4pt);
  \node[font=\scriptsize\bfseries,above] at ({\mx{0}},0.72) {Vega};

  \fill[milcNegro!55] ({\mx{4}},0.5) circle (1.7pt);
  \fill[milcNegro!45] ({\mx{6}},0.5) circle (1pt);

  % ── Zonas: lo que se pierde por el alumbrado ─────────────────────────
  \draw[decorate,decoration={brace,amplitude=4pt},milcGradeDark,line width=.7pt]
    ({\mx{6}},-0.62) -- ({\mx{4}},-0.62);
  \node[font=\scriptsize\bfseries,milcGradeDark,align=center,below]
    at ({\mx{5}},-0.72) {estas estrellas te las\\borra el alumbrado};

  \draw[dashed,milcTurquesa,line width=.8pt] ({\mx{4}},-0.45) -- ({\mx{4}},1.35);
  \node[font=\scriptsize\bfseries,milcTurquesa,align=center,above]
    at ({\mx{4}},1.3) {Cartago\\centro};

  \draw[dashed,milcVerde,line width=.8pt] ({\mx{6}},-0.45) -- ({\mx{6}},1.35);
  \node[font=\scriptsize\bfseries,milcVerde,align=center,above]
    at ({\mx{6}},1.3) {vereda\\oscura};

  % (La síntesis va en el pie de figura vía \guiaDiagrama, para no repetirla
  %  dentro del propio diagrama.)

\end{tikzpicture}
